\documentclass[11pt]{article}
\usepackage[T1]{fontenc}
\usepackage[french]{babel}
\usepackage{amsmath,amssymb,mathtools,amscd,esint,mathrsfs,bbm,bm,cancel,stmaryrd}
\usepackage{tikz-cd}
\usepackage[all]{xy}
\usepackage{physics}
\usepackage{siunitx}
\usepackage[version=4]{mhchem}
\usepackage{chemfig}
\usepackage{tensor}
\usepackage{ytableau}
\usepackage{nicematrix}
\usepackage{hyperref}

\DeclarePairedDelimiter\valabs{\lvert}{\rvert}
\newcommand{\R}{\mathbb{R}}
\newcommand{\norme}[2][2]{\left\lVert #2 \right\rVert_{#1}}
\makeatletter
\def\produit#1#2{\left\langle #1 \,\middle|\, #2 \right\rangle}
\makeatother

\begin{document}

\section[Diagrammes]{Diagrammes $\mathscr{C}$ et \texorpdfstring{$\infty$}{infini}-catégories}\label{sec:diag}
Un carré commutatif (tikz-cd) :
\[
\begin{tikzcd}
A \arrow[r, "f"] \arrow[d, "g"'] & B \arrow[d, "h"] \\
C \arrow[r, "k"'] \arrow[ur, dashed, "\exists!"] & D
\end{tikzcd}
\qquad
\begin{CD}
X @>\varphi>> Y \\
@VVV @AAA \\
Z @= Z
\end{CD}
\qquad
\xymatrix{ E \ar[r]^{p} & F \ar@{->>}[d] \\ & G }
\]

\section{Analyse}
\begin{equation}\label{eq:Navier–Stokes_2}
\oiint_{\partial V} \bm{u}\cdot\dd{\bm{S}} = \iiint_V \div{\bm{u}} \dd{V},
\qquad \pdv[2]{f}{x} + \dv{g}{t} = \comm{A}{B} + \bra{\psi}\ket{\phi}.
\end{equation}
Une fraction continue :
\[
x = a_0 + \cfrac{1}{a_1 + \cfrac{1}{a_2 + \cfrac{1}{a_3 + \cfrac{1}{\ddots}}}}
\quad
\sum_{\substack{0 \le i \le m \\ 0 < j < n}} P(i,j)
\quad
\sideset{_*^*}{'}\sum_{k} a_k
\quad
\operatorname*{arg\,max}_{x \in \R} \valabs*{\frac{f(x)}{g(x)}}
\]
\[
\overbrace{a + \underbrace{b + c}_{\text{deux}} + d}^{\text{trois}}
\overset{!}{=} \underset{n\to\infty}{\lim} u_n
\coloneqq \prescript{14}{6}{\mathrm{C}}
\xRightarrow[\text{bas}]{\text{haut}} \cancelto{0}{x} + \bcancel{y} + \xcancel{z}
\]
\[
\tensor{R}{^a_b_c^d} \qquad \mathbbm{1}_A \qquad \llbracket 1, n \rrbracket \qquad \norme[\infty]{f}
\qquad \produit{\varphi}{\psi}
\qquad
\begin{dcases}
\frac{1}{x} & \text{si } x > 0 \\
\int_0^1 t^2 \, \mathrm{d}t & \text{sinon}
\end{dcases}
\]
\[
\begin{pNiceMatrix}[first-row,last-col]
C_1 & \Cdots & C_n & \\
a_{11} & \Cdots & a_{1n} & L_1 \\
\Vdots & \Ddots & \Vdots & \Vdots \\
a_{n1} & \Cdots & a_{nn} & L_n
\end{pNiceMatrix}
\qquad
\left(\begin{array}{c|cc}
1 & 0 & 0 \\ \hline
0 & \alpha & \beta
\end{array}\right)
\]

\section{Physique et chimie}
$g = \SI{9.81}{\metre\per\second\squared}$, $e = \num{1.602e-19}$ C, 50\% des cas.
\ce{2H2 + O2 -> 2H2O}, \ce{^{14}_{6}C}, \chemfig{*6(=-=-=-)}, \ydiagram{3,2,1}.
Retour à la ligne puis commentaire : a\\% ceci est un vrai commentaire
Voir l'équation~\eqref{eq:Navier–Stokes_2}, la section~\ref{sec:diag} et~\cite[th.~2]{knuth, lamport}.

\begin{verbatim}
\section{Faux titre dans du verbatim}
% pas un commentaire ici
\end{verbatim}
\verb|\section{aussi faux}|

\section{Théorie des groupes}\label{sec:groupes}
Pour $\sigma \in \mathfrak{S}_n$, la signature vérifie
\[
\varepsilon(\sigma) = \prod_{1 \le i < j \le n} \frac{\sigma(j) - \sigma(i)}{j - i} \in \{-1, 1\}.
\]


\bibliographystyle{plain}
\bibliography{refs}
\end{document}
